\documentclass[11pt,a4paper]{article}
\usepackage{gtu-el-yazisi}
\begin{document}
\notbilgi{Görüntü İşleme}{Morphological Image Processing (devam)}{Ders-11.pdf}
\textit{Not: Orijinal dosyadaki basılı sayfalar (s. 2, 7) dijitalleştirilmedi; yalnızca el yazısı sayfalar aktarıldı.}

\sayfa{1}
\section{Conditional dilation}
\[ A \oplus_C B = (A \oplus B) \cap C \]

(c)

\[ A = C \circ D \qquad \text{(opening)} \qquad \leftarrow \text{D büyük olsun} \]

\elcizimi[0.85]{sekiller/Ders-11-s1-1.png}[1]

Connected component bulmak için
\[ X_k = (X_{k-1} \oplus B) \cap A \]
\[ k = 1,2,\dots \]
$X_0$ ihtiyaç var, objenin üzerindeki herhangi bir nokta

\sekilnotu[1]{Bu sayfa basılı materyal içeriyor (ör. ders slaytı); dijitalleştirilmedi.}
$\rightarrow$ 8 komşuluk

dilation and intersection

$X_0$ bulmak için etiketlemenin herhangi bir noktası

\sayfa{3}
$x$ bir değişken
\[ f(x) \]
\[ f_x(z) = f(z-x) \qquad \text{translation by x} \]
\[ (f+y)(z) = f(z)+y \qquad \text{offset by y} \]
\[ f_x + y \]

\elcizimi[0.7]{sekiller/Ders-11-s3-1.png}[3]

\[ D[g] \qquad D[f] \]

$g$ and $f$ are signals with domain $D[g]$ and $D[f]$ we say that $g$ is beneath of (altında) $f$, if

1) $D[g] \subseteq D[f]$

2) $g(x) \le f(x)$ for any $x \in D[g]$

Then we write $g \ll f$

\elcizimi[0.5]{sekiller/Ders-11-s3-2.png}[3]

$\rightarrow g \ll f$

\elcizimi[0.8]{sekiller/Ders-11-s3-3.png}[3]

\sayfa{4}
\[ (f \wedge g)(x) = \min\{f(x), g(x)\} \quad \text{for } x \in D[f] \cap D[g] \]
\[ (f \vee g)(x) = \begin{cases}
\max\{f(x), g(x)\} & \text{for } x \in D[g] \cap D[f]\\
f(x) & \text{for } x \in D[f] \text{ and } x \notin D[g]\\
g(x) & \text{for } x \in D[g] \text{ and } x \notin D[f]\\
\text{not defined} & \text{for } x \notin D[g] \text{ and } x \notin D[f]
\end{cases} \]

\elcizimi[0.85]{sekiller/Ders-11-s4-1.png}[4]

\section{Reflection of h}
\[ \hat{h}(x) = -h(-x) \]

\elcizimi[0.35]{sekiller/Ders-11-s4-2.png}[4]

\[ (f \ominus g)(x) = \max\{y : g_x + y \ll f\} \]

\elcizimi[0.7]{sekiller/Ders-11-s4-3.png}[4]

\sayfa{5}
\elcizimi[0.3]{sekiller/Ders-11-s5-1.png}[5]

\begin{align*}
(f \oplus g)(x) &= \min\{y : (\hat{g})_x + y \gg f\}\\
&= \overset{\text{max}}{\bigvee}\{f_x + g(x) : x \in D[g]\}
\end{align*}

\elcizimi[0.9]{sekiller/Ders-11-s5-2.png}[5]

\begin{align*}
f \circ g &= (f \ominus g) \oplus g\\
f \bullet g &= (f \oplus g) \ominus g\\
(f \circ g) \circ g &= f \circ g
\end{align*}

\begin{align*}
\text{Grad}(f) &= f \oplus g - f \ominus g\\
\text{Top-Hat}(f) &= f - f \circ g
\end{align*}

\elcizimi[0.5]{sekiller/Ders-11-s5-3.png}[5]

opening olayı background estimation olarak yarar

$f \circ g$ background

$\swarrow$ arkası karanlık üstü beyaz şeyler için

\sekilnotu[5]{Bu sayfa basılı materyal içeriyor (ör. ders slaytı); dijitalleştirilmedi.}

\sayfa{6}
\[ \text{Bottom Hat}(f) = f - f \bullet g\footnotemark \]
\footnotetext{Claude notu: muhtemelen $f \bullet g - f$ olmalı.}

$\downarrow$ Gradient
\[ A \oplus B - A \ominus B \]

Morphological laplacian
\[ \frac{1}{2}(A \oplus B + A \ominus B - 2A) \]
\[ \frac{1}{2}\Big(\underbrace{(A \oplus B - A)}_{\text{türev}} - \underbrace{(A - A \ominus B)}_{\text{türev}}\Big) \]

$\downarrow$ smoothing
\[ (A \circ B) \bullet B \]

\elcizimi[0.45]{sekiller/Ders-11-s6-1.png}[6]

minimum
\[ \frac{A - A \ominus B}{A \oplus B - A \ominus B} \;\Big\}\; \text{normalizasyon gibi} \]
$\searrow$ B biraz büyük seçilmeli

contrast stretching local

erosion karanlık alanları büyütüyor

\sekilnotu[6]{Bu sayfa basılı materyal içeriyor (ör. ders slaytı); dijitalleştirilmedi.}

\end{document}
